\section{Related Work}
\label{sec:related}

The task addressed in this paper lies at the intersection of fact-checking and information retrieval. A natural starting point is lexical retrieval where probabilistic ranking models such as BM25 remain strong baselines thanks to their efficiency, interpretability and robustness in candidate generation \citep{RobertsonZaragoza2009}. In the setting of scientific claim source retrieval, lexical matching is especially useful when a post mentions recognizable entities or terminology that also appears in the title, abstract, venue or author list of the referenced paper.

Neural retrieval methods have been proposed to mitigate the lexical mismatch between queries and documents by learning dense semantic representations. More broadly, neural information retrieval has shifted attention from exact token overlap to representation learning and learned matching functions \citep{MarchesinEtAl2020}. In particular, dense dual-encoder architectures such as Dense Passage Retrieval (DPR) demonstrated that learned query and document embeddings can be highly effective for first-stage retrieval \citep{KarpukhinEtAl2020}. This line of work is especially relevant to our task because social media posts often paraphrase scientific findings instead of reusing the exact wording found in the original publication. BGE-M3 unified dense, sparse, and multi-vector retrieval in a single multilingual model \citep{ChenEtAl2024}, 
enabling hybrid scoring that combines embedding similarity with lexical token weights a combination particularly well-suited to cross-lingual settings where queries and documents share few surface tokens.

Vocabulary mismatch between short, informal queries and longer scientific documents can be further reduced through query expansion. Classical pseudo-relevance feedback methods such as RM3 \citep{LavrenkoEtAl2001} re-weight query terms using top-ranked documents from an initial retrieval pass. Neural extensions of this paradigm apply dense representations to the same idea \citep{WangEtAl2021PRF}. In our pipeline, we incorporate embedding-based query expansion to enrich query representations before scoring against the corpus.

Between initial retrieval and final ranking, many modern systems adopt multi-stage pipelines that trade off efficiency and effectiveness \citep{MarchesinEtAl2020}. Late-interaction models such as ColBERT preserve fine-grained token-level matching while remaining substantially more scalable than fully pairwise neural scoring over the whole collection \citep{KhattabZaharia2020}. Likewise, transformer-based re-rankers have become a common second-stage component for refining an initial candidate set with stronger query-document interactions \citep{NogueiraEtAl2020}. These approaches motivate retrieval architectures that first generate a manageable set of candidates and then apply more expensive relevance estimation only where it is most useful.

Recent work has argued that reranker fine-tuning benefits from hard negatives because they force the model to learn finer-grained relevance 
boundaries than easy or random negatives provide \citep{MeghwaniEtAl2025HardNegatives}. A practical instantiation of this strategy trains the reranker on negatives mined from the same retrieval model whose errors the reranker is meant to correct, a form of bootstrap 
self-improvement. Adapting large language models to retrieval tasks is further facilitated by parameter-efficient fine-tuning with LoRA \citep{HuEtAl2022}, which has been shown to match or approach full fine-tuning quality for passage reranking while dramatically reducing GPU memory requirements \citep{nijasure2025relevance}. This perspective is consistent with recent practical guidance for training cross-encoder rerankers with Sentence Transformers \citep{Aarsen2025TrainReranker}. In that setting, fast embedding-based retrieval is used to build a candidate pool, while the cross-encoder jointly scores query-document pairs and benefits from negatives mined from nearby embedding neighborhoods rather than from trivial mismatches. 

Our work is also related to the broader literature on automatic claim verification and more specifically to the CheckThat! shared tasks, which have addressed the detection, verification, and analysis of claims in social media across several editions \citep{ElsayedEtAl2019,BarronCedenoEtAl2020,NakovEtAl2021}. Recent participants in the same task line have  independently converged on hybrid retrieval plus reranking as the dominant pipeline configuration \citep{Sager2025, Schofield2025}, confirming the high-level design choices taken.